\chapter{Design}

\section{System Architecture}
The Intelligent Security Compliance Management System (ISCMS) uses an N-tier architecture. Separating the frontend, API, background workers, and database allows for scaling and debugging of individual components without affecting the entire system. Figure \ref{fig:tier_architecture} shows the four main layers.

\subsection{Client and Presentation Tier}
The frontend is a React 18 Single Page Application (SPA). Moving the rendering logic to the browser reduces server load and lets analysts review compliance alerts without full page reloads.

\subsection{Application and Middleware Tier}
The backend uses the Django REST Framework. It is the primary gateway, verifying JWTs and enforcing role-based access before processing queries. Database traffic is routed through Django to maintain a consistent security boundary.

\subsection{Asynchronous Processing Tier}
Retrieving telemetry from the Wazuh SIEM is an I/O-intensive process. Executing these tasks on the main Django thread would cause API latency and timeouts. To prevent this, the system offloads I/O tasks to background Celery workers using Redis as the message broker.

\subsection{Data Persistence Tier}
PostgreSQL 18 is the primary database for its ACID guarantees. The schema uses relational tables for mapping frameworks to controls, while the JSONField type stores unstructured SIEM logs.

\begin{figure}[H]
    \centering
    \resizebox{0.85\textwidth}{!}{
    \begin{tikzpicture}[
        tier/.style={rectangle, draw=blue!80!black, thick, fill=blue!5, minimum width=8.5cm, minimum height=1.6cm, align=center, font=\bfseries},
        arrow/.style={-Stealth, thick, draw=black!80},
        label/.style={font=\scriptsize, inner sep=2pt}
    ]
        % Tiers with increased separation to resolve congestion
        \node[tier] (client) at (0, 8.5) {Client \& Presentation Tier\\(React SPA)};
        \node[tier] (app) at (0, 5.0) {Application \& Middleware Tier\\(Django REST Framework)};
        
        \node[tier, minimum width=4cm] (async) at (-3.2, 1.5) {Async Processing\\(Celery \& Redis)};
        \node[tier, minimum width=4cm] (data) at (3.2, 1.5) {Data Persistence\\(PostgreSQL)};
        
        \node[tier, fill=gray!10, draw=gray!80!black] (wazuh) at (-3.2, -1.5) {External Security\\(Wazuh SIEM)};

        % Client <-> App Connections
        \draw[arrow] ([xshift=-2cm]client.south) -- ([xshift=-2cm]app.north) 
            node[midway, left=8pt, label] {HTTPS / JSON};
        \draw[arrow] ([xshift=2cm]app.north) -- ([xshift=2cm]client.south) 
            node[midway, right=8pt, label] {API Responses};

        % App -> Async / Data Connections (Widened entry points)
        \draw[arrow] ([xshift=-2cm]app.south) -- ([xshift=0.8cm]async.north) 
            node[midway, left=10pt, label] {Enqueue Tasks};
        \draw[arrow] ([xshift=2cm]app.south) -- ([xshift=-0.8cm]data.north) 
            node[midway, right=10pt, label] {ORM Queries};

        % Async -> Wazuh (Increased vertical gap)
        \draw[arrow] (async.south) -- (wazuh.north) 
            node[midway, left=8pt, label] {REST API Polling};
        
        % Async -> Data (Increased horizontal gap)
        \draw[arrow] (async.east) -- (data.west) 
            node[midway, above=6pt, label] {State Commits};
    \end{tikzpicture}
    }
    \caption{Tier-wise System Architecture of the ISCMS Platform}
    \label{fig:tier_architecture}
\end{figure}

\section{Design Constraints}
The design of the ISCMS involved addressing several technical and environmental constraints.

\subsection{Technical Constraints}
For security, the database and Redis broker operate inside a private Docker bridge network. Their ports are not exposed to the host machine. Access to data requires requests through the Django API container.

The integration with the Wazuh API introduced rate limits and network latency. Celery workers are configured to batch and throttle polling requests. Consequently, the dashboard displays the compliance state from the last successful synchronization rather than a real-time feed.

On the frontend, rendering large compliance grids can impact browser performance. Loading thousands of DOM nodes simultaneously can block the main thread. React virtualization addresses this by rendering only the visible elements.

\subsection{Infrastructure and Operational Constraints}
External factors, such as restricted access to university facilities, required the transition to developing and testing the architecture in a home-lab environment.

Running a full Wazuh SIEM manager locally is resource-intensive and can impact system stability. To manage this, a distributed configuration was used. One system was dedicated to running Wazuh, while another was used for development and as a monitored endpoint. Git was used to maintain consistency across different development environments while the primary infrastructure remained on a dedicated server.

\section{Design Methodology}
The platform was built using an iterative prototyping approach. The system was divided into distinct modules to allow for independent testing before integration.

The project was divided into four main phases:

\begin{itemize}
    \item \textbf{Phase 1 (Prototyping):} The initial phase focused on verifying the core concept. This included setting up a basic Django backend and testing the Wazuh API to confirm successful retrieval of Security Configuration Assessment (SCA) telemetry.
    
    \item \textbf{Phase 2 (Core Architecture):} The second phase involved designing the PostgreSQL database schema to link frameworks to rules. JWT authentication and role-based routing were also implemented to secure the API.
    
    \item \textbf{Phase 3 (Asynchronous Telemetry):} Celery and Redis were added to handle heavy SIEM data ingestion. This established a background worker fleet for data processing and deduplication.
    
    \item \textbf{Phase 4 (Frontend \& Visualization):} The final phase involved building the React frontend. This included rendering data grids, implementing state management, and configuring the server-side PDF generator for audit reports.
\end{itemize}

\section{High Level Design}
The system architecture is represented through several views that cover software components, physical deployment, and operational sequences.

\subsection{Component Diagram}
Figure \ref{fig:component_diagram} shows the interactions between system modules. The React frontend sends HTTPS requests to the Django API Gateway. The IAM module verifies the JWT, and the Compliance Engine then translates the request into a database query. A separate Wazuh Polling Service retrieves alerts and stores them in PostgreSQL.

\begin{figure}[H]
    \centering
    \resizebox{0.9\textwidth}{!}{
    \begin{tikzpicture}[
        comp/.style={rectangle, draw=black!80, thick, fill=white, minimum width=4.8cm, minimum height=2cm, text width=3.8cm, align=center, font=\small\bfseries, rounded corners=2pt,
            path picture={
                \draw[black!80, thick, fill=gray!10] ([xshift=-8mm,yshift=-2mm]path picture bounding box.north east) rectangle ++(6mm,-6mm);
                \draw[black!80, thick, fill=white] ([xshift=-10mm,yshift=-3mm]path picture bounding box.north east) rectangle ++(4mm,-1.5mm);
                \draw[black!80, thick, fill=white] ([xshift=-10mm,yshift=-6.5mm]path picture bounding box.north east) rectangle ++(4mm,-1.5mm);
            }
        },
        extcomp/.style={comp, fill=gray!10},
        arrow/.style={-Stealth, thick, dashed, draw=black!70},
        label/.style={font=\scriptsize, fill=white, inner sep=2pt}
    ]
        % System Boundary (Shifted up slightly for title clearance)
        \node[draw=blue!80!black, thick, minimum width=19.5cm, minimum height=8.2cm, fill=blue!2, rounded corners=5pt, dashed] (system) at (3, 0.85) {};
        \node[anchor=north west, xshift=2mm, yshift=-1mm, font=\bfseries\large, text=blue!80!black] at (system.north west) {ISCMS Application Subsystem};

        % ROW 0: External Top
        \node[extcomp] (wazuh) at (9.5, 7.2) {External Wazuh\\Manager API};

        % ROW 1: Internal Top
        \node[comp] (ui) at (-3.5, 3) {React Frontend\\UI Component};
        \node[comp] (gateway) at (3, 3) {Django DRF\\API Gateway};
        
        % ROW 2: Internal Bottom
        \node[comp] (auth) at (-3.5, -1.5) {IAM \& Session\\Manager};
        \node[comp] (engine) at (3, -1.5) {Compliance\\Mapping Engine};
        \node[comp] (poll) at (9.5, -1.5) {Wazuh Polling\\Celery Service};

        % ROW 3: External Bottom
        \node[extcomp] (db) at (3, -6) {PostgreSQL\\Database};
        \node[extcomp] (redis) at (9.5, -6) {Redis Message\\Broker};

        % Dependencies
        \draw[arrow] (ui) -- (gateway) node[midway, below=5pt, label] {REST / JSON};
        \draw[arrow] (gateway) -- (auth) node[midway, above=5pt, label, sloped] {JWT Validation};
        \draw[arrow] (gateway) -- (engine) node[midway, right=5pt, label] {Data Requests};
        \draw[arrow] (gateway) -- (wazuh) node[midway, above=5pt, label, sloped] {Trigger Sync};
        \draw[arrow] (wazuh) -- (poll) node[midway, right=5pt, label] {HTTPS Telemetry};
        
        \draw[arrow] (engine) -- (db) node[midway, right=5pt, label] {ORM Queries};
        \draw[arrow] (poll) -- (db) node[midway, above=5pt, label, sloped] {State Commits};
        \draw[arrow] (poll) -- (redis) node[midway, right=5pt, label] {Task Queues};
    \end{tikzpicture}
    }
    \caption{Conceptual View: UML Component Diagram}
    \label{fig:component_diagram}
\end{figure}

\subsection{Deployment Diagram}
Figure \ref{fig:deployment_diagram} outlines the Docker container environment. The Django backend, Celery workers, Redis, and PostgreSQL are deployed on the same host but isolated within a private bridge network. This keeps the database inaccessible from external networks.

\begin{figure}[H]
    \centering
    \resizebox{0.95\textwidth}{!}{
    \begin{tikzpicture}[
        device/.style={
            rectangle, draw=black!90, thick, fill=white,
            minimum width=4cm, minimum height=2.5cm,
            append after command={
                \pgfextra{
                    \filldraw[black!90, thick, fill=gray!15] (\tikzlastnode.north west) -- ++(0.3, 0.4) -- ($(\tikzlastnode.north east)+(0.3,0.4)$) -- (\tikzlastnode.north east) -- cycle;
                    \filldraw[black!90, thick, fill=gray!15] (\tikzlastnode.north east) -- ($(\tikzlastnode.north east)+(0.3,0.4)$) -- ($(\tikzlastnode.south east)+(0.3,0.4)$) -- (\tikzlastnode.south east) -- cycle;
                }
            }
        },
        artifact/.style={
            rectangle, draw=black!80, thick, fill=blue!5,
            minimum width=3.4cm, minimum height=1cm,
            align=center, font=\scriptsize\bfseries, rounded corners=2pt
        },
        conn/.style={
            thick, draw=black!80, Stealth-Stealth
        },
        label/.style={
            font=\scriptsize, fill=white, inner sep=2pt
        }
    ]
        % ----------------------------------------------------
        % CENTER: Docker Host Server
        % ----------------------------------------------------
        \node[device, minimum width=9cm, minimum height=6cm] (server) at (0, 0) {};
        \node[anchor=north, font=\small\bfseries, yshift=-2mm, align=center] at (server.north) {\textless\textless device\textgreater\textgreater\\Docker Host Server (Ubuntu)};
        
        \node[artifact] (django) at (-2.4, 1) {\textless\textless container\textgreater\textgreater\\Django API Gateway};
        \node[artifact] (celery) at (2.4, 1) {\textless\textless container\textgreater\textgreater\\Celery Async Worker};
        \node[artifact] (postgres) at (-2.4, -1.5) {\textless\textless container\textgreater\textgreater\\PostgreSQL Database};
        \node[artifact] (redis) at (2.4, -1.5) {\textless\textless container\textgreater\textgreater\\Redis Message Broker};

        % Internal Bridge Network Connections
        \draw[dashed, -Stealth, thick, draw=black!70] (django.south) -- (postgres.north) node[midway, fill=white, font=\tiny] {TCP 5432};
        \draw[dashed, -Stealth, thick, draw=black!70] (celery.south) -- (redis.north) node[midway, fill=white, font=\tiny] {TCP 6379};
        \draw[dashed, Stealth-Stealth, thick, draw=black!70] (django.east) -- (celery.west) node[midway, below=3pt, fill=white, font=\tiny] {Task IPC};

        % ----------------------------------------------------
        % LEFT: Client Workstation
        % ----------------------------------------------------
        \node[device, minimum width=4cm, minimum height=3cm] (client) at (-8.5, 1.5) {};
        \node[anchor=north, font=\small\bfseries, yshift=-2mm, align=center] at (client.north) {\textless\textless device\textgreater\textgreater\\Client Workstation};
        \node[artifact] (react) at (-8.5, 1) {\textless\textless artifact\textgreater\textgreater\\React SPA Browser};

        % ----------------------------------------------------
        % RIGHT: Wazuh SIEM Node
        % ----------------------------------------------------
        \node[device, minimum width=4cm, minimum height=3cm] (wazuh) at (8.5, 1.5) {};
        \node[anchor=north, font=\small\bfseries, yshift=-2mm, align=center] at (wazuh.north) {\textless\textless device\textgreater\textgreater\\Wazuh SIEM Node};
        \node[artifact] (sca) at (8.5, 1) {\textless\textless artifact\textgreater\textgreater\\SCA Telemetry API};

        % ----------------------------------------------------
        % EXTERNAL NETWORK CONNECTIONS (Perfectly Horizontal at Y=1)
        % ----------------------------------------------------
        \draw[conn] (react.east) -- (django.west) node[midway, below=4pt, label] {HTTPS / TLS 1.3};
        \draw[conn, -Stealth] (celery.east) -- (sca.west) node[midway, below=4pt, label] {REST API Polling};

    \end{tikzpicture}
    }
    \caption{Physical View: UML Deployment Diagram}
    \label{fig:deployment_diagram}
\end{figure}

\subsection{System Activity Diagrams}
The activity diagrams below show the control flow for core platform features. Standard UML semantics, such as decision nodes and synchronization forks, are used to represent the system logic.

\subsubsection{Authentication and IAM Routing}
Figure \ref{fig:activity_auth} shows the login process. Django verifies credentials against PostgreSQL. If the verification fails, the system logs the event and denies the request. Upon successful authentication, the system issues a JWT and evaluates the user's role. Administrators are directed to configuration settings, while Auditors are restricted to read-only dashboards.

\begin{figure}[H]
    \centering
    \resizebox{0.75\textwidth}{!}{
    \begin{tikzpicture}[
        node distance=1.5cm,
        startstop/.style={circle, fill=black, minimum size=0.4cm},
        end/.style={circle, draw=black, thick, minimum size=0.4cm, path picture={\fill[black] (path picture bounding box.center) circle (0.15cm);}},
        activity/.style={rectangle, draw=black!80, thick, fill=blue!5, minimum width=3.5cm, minimum height=1cm, align=center, font=\scriptsize\bfseries, rounded corners=2pt, text width=3.2cm},
        decision/.style={diamond, draw=black!80, thick, fill=yellow!10, minimum width=2.5cm, minimum height=2.5cm, align=center, font=\scriptsize\bfseries, aspect=1.5, text width=1.8cm},
        arrow/.style={-Stealth, thick, draw=black!90},
        label/.style={font=\scriptsize, fill=white, inner sep=1pt}
    ]
        \node[startstop] (start) at (0,0) {};
        \node[activity] (submit) at (0,-1.5) {Submit Credentials via HTTPS};
        \node[activity] (verify) at (0,-3) {Verify Hash in PostgreSQL};
        \node[decision] (isvalid) at (0,-5.5) {Credentials Valid?};
        \node[activity] (reject) at (5,-5.5) {Return HTTP 401 \& Log Event};
        \node[activity] (jwt) at (0,-8.5) {Generate JWT Access \& Refresh Pair};
        \node[decision] (role) at (0,-11.5) {Evaluate RBAC Profile};
        \node[activity] (admin) at (-3,-14) {Route to System Configuration};
        \node[activity] (auditor) at (3,-14) {Route to Read-Only Dashboard};
        \node[end] (end1) at (5,-7.5) {};
        \node[end] (end2) at (0,-16) {};

        \draw[arrow] (start) -- (submit);
        \draw[arrow] (submit) -- (verify);
        \draw[arrow] (verify) -- (isvalid);
        \draw[arrow] (isvalid) -- (reject) node[midway, above, label] {No};
        \draw[arrow] (reject) -- (end1);
        \draw[arrow] (isvalid) -- (jwt) node[midway, left, label] {Yes};
        \draw[arrow] (jwt) -- (role);
        \draw[arrow] (role) -| (admin) node[pos=0.55, above, label] {Superadmin};
        \draw[arrow] (role) -| (auditor) node[pos=0.55, above, label] {Auditor};
        \draw[arrow] (admin) |- (end2);
        \draw[arrow] (auditor) |- (end2);
    \end{tikzpicture}
    }
    \caption{Activity Diagram: Authentication and IAM Routing}
    \label{fig:activity_auth}
\end{figure}

\subsubsection{Asynchronous Telemetry Ingestion}
Figure \ref{fig:activity_sync} shows the asynchronous ingestion process. When a synchronization task is triggered, the main thread returns an HTTP 202 status to the client. Simultaneously, the Celery worker polls the Wazuh API. If the SIEM is unreachable, the system uses exponential backoff. Once telemetry is retrieved, the system parses the JSON data and updates the database.

\begin{figure}[H]
    \centering
    \resizebox{0.8\textwidth}{!}{
    \begin{tikzpicture}[
        node distance=1.5cm,
        startstop/.style={circle, fill=black, minimum size=0.4cm},
        end/.style={circle, draw=black, thick, minimum size=0.4cm, path picture={\fill[black] (path picture bounding box.center) circle (0.15cm);}},
        activity/.style={rectangle, draw=black!80, thick, fill=blue!5, minimum width=3.5cm, minimum height=1cm, align=center, font=\scriptsize\bfseries, rounded corners=2pt, text width=3.2cm},
        decision/.style={diamond, draw=black!80, thick, fill=yellow!10, minimum width=2cm, minimum height=2cm, align=center, font=\scriptsize\bfseries, aspect=1.5, text width=1.5cm},
        fork/.style={rectangle, draw=black, fill=black, minimum width=8cm, minimum height=0.15cm},
        arrow/.style={-Stealth, thick, draw=black!90},
        label/.style={font=\scriptsize, fill=white, inner sep=1pt}
    ]
        \node[startstop] (start) at (0,0) {};
        \node[activity] (trigger) at (0,-1.5) {Trigger Sync via Cron or API};
        \node[fork] (fork1) at (0,-3) {};
        
        \node[activity] (ui) at (-3,-4.5) {UI Thread: Display Syncing Indicator};
        \node[activity] (celery) at (3,-4.5) {Worker: Poll Wazuh SCA API};
        \node[decision] (api) at (3,-7) {API OK?};
        \node[activity] (retry) at (7,-7) {Exponential Backoff \& Log};
        \node[activity] (parse) at (3,-9.5) {Parse JSON Telemetry Array};

        \node[fork] (join1) at (0,-11.5) {};
        \node[activity] (commit) at (0,-13) {Execute Bulk PostgreSQL Commit};
        \node[end] (end) at (0,-15) {};

        \draw[arrow] (start) -- (trigger);
        \draw[arrow] (trigger) -- (fork1);
        \draw[arrow] (fork1) -| (ui);
        \draw[arrow] (fork1) -| (celery);
        \draw[arrow] (celery) -- (api);
        \draw[arrow] (api) -- (retry) node[midway, above, label] {No};
        \draw[arrow] (retry) |- (fork1);
        \draw[arrow] (api) -- (parse) node[midway, left, label] {Yes};
        
        \draw[arrow] (ui) |- (join1);
        \draw[arrow] (parse) |- (join1);
        \draw[arrow] (join1) -- (commit);
        \draw[arrow] (commit) -- (end);
    \end{tikzpicture}
    }
    \caption{Activity Diagram: Parallel Asynchronous Telemetry Ingestion}
    \label{fig:activity_sync}
\end{figure}

\subsubsection{Compliance State Evaluation}
The mapping of technical alerts to policies like ISO 27001 is shown in Figure \ref{fig:activity_eval}. The system uses a worst-case priority algorithm; if a technical check fails on a single node, the corresponding governance control is marked as non-compliant. This evaluation continues until all policies are processed.

\begin{figure}[H]
    \centering
    \resizebox{0.75\textwidth}{!}{
    \begin{tikzpicture}[
        node distance=1.5cm,
        startstop/.style={circle, fill=black, minimum size=0.4cm},
        end/.style={circle, draw=black, thick, minimum size=0.4cm, path picture={\fill[black] (path picture bounding box.center) circle (0.15cm);}},
        activity/.style={rectangle, draw=black!80, thick, fill=blue!5, minimum width=3.5cm, minimum height=1cm, align=center, font=\scriptsize\bfseries, rounded corners=2pt, text width=3.2cm},
        decision/.style={diamond, draw=black!80, thick, fill=yellow!10, minimum width=2.5cm, minimum height=2.5cm, align=center, font=\scriptsize\bfseries, aspect=1.5, text width=2cm},
        arrow/.style={-Stealth, thick, draw=black!90},
        label/.style={font=\scriptsize, fill=white, inner sep=1pt}
    ]
        \node[startstop] (start) at (0,0) {};
        \node[activity] (load) at (0,-1.5) {Load ISO 27001 Control Matrix};
        \node[activity] (fetch) at (0,-3) {Fetch Current Node Telemetry};
        \node[decision] (check) at (0,-5.5) {Any Checks Failed?};
        \node[activity] (fail) at (-3,-8) {Mark Control Non-Compliant};
        \node[activity] (pass) at (3,-8) {Mark Control Compliant};
        \node[decision] (more) at (0,-10.5) {More Controls?};
        \node[activity] (render) at (0,-13) {Update Aggregated Compliance Score};
        \node[end] (end) at (0,-15) {};

        \draw[arrow] (start) -- (load);
        \draw[arrow] (load) -- (fetch);
        \draw[arrow] (fetch) -- (check);
        \draw[arrow] (check) -| (fail) node[pos=0.45, above, label] {Yes (Worst-Case)};
        \draw[arrow] (check) -| (pass) node[pos=0.45, above, label] {No};
        \draw[arrow] (fail) |- (more);
        \draw[arrow] (pass) |- (more);
        \draw[arrow] (more) -- (render) node[midway, left, label] {No};
        \draw[arrow] (more) -- ++(5,0) |- (fetch) node[pos=0.25, right, label] {Yes};
        \draw[arrow] (render) -- (end);
    \end{tikzpicture}
    }
    \caption{Activity Diagram: Iterative Compliance Evaluation and Worst-Case Priority}
    \label{fig:activity_eval}
\end{figure}

\subsubsection{Audit Report Generation}
The audit report generation process is automated, as shown in Figure \ref{fig:activity_report}. When a user selects a framework, the backend aggregates relevant database records. This data is injected into an HTML template, and a PDF is generated on the server for delivery to the client.

\begin{figure}[H]
    \centering
    \resizebox{0.6\textwidth}{!}{
    \begin{tikzpicture}[
        node distance=1.5cm,
        startstop/.style={circle, fill=black, minimum size=0.4cm},
        end/.style={circle, draw=black, thick, minimum size=0.4cm, path picture={\fill[black] (path picture bounding box.center) circle (0.15cm);}},
        activity/.style={rectangle, draw=black!80, thick, fill=blue!5, minimum width=4cm, minimum height=1cm, align=center, font=\scriptsize\bfseries, rounded corners=2pt, text width=3.8cm},
        decision/.style={diamond, draw=black!80, thick, fill=yellow!10, minimum width=2.5cm, minimum height=2.5cm, align=center, font=\scriptsize\bfseries, aspect=1.5, text width=2cm},
        arrow/.style={-Stealth, thick, draw=black!90},
        label/.style={font=\scriptsize, fill=white, inner sep=1pt}
    ]
        \node[startstop] (start) at (0,0) {};
        \node[activity] (request) at (0,-1.5) {Initialize Report Request (Framework ID)};
        \node[activity] (query) at (0,-3) {Aggregate Historical Compliance Data};
        \node[decision] (data) at (0,-5.5) {Data Exists?};
        \node[activity] (error) at (5,-5.5) {Abort \& Alert Auditor};
        \node[activity] (template) at (0,-8.5) {Inject Data into Reporting Template};
        \node[activity] (pdf) at (0,-10) {Execute PDF Generation Engine};
        \node[activity] (download) at (0,-11.5) {Transmit Binary Payload to Client};
        \node[end] (end1) at (5,-7.5) {};
        \node[end] (end2) at (0,-13.5) {};

        \draw[arrow] (start) -- (request);
        \draw[arrow] (request) -- (query);
        \draw[arrow] (query) -- (data);
        \draw[arrow] (data) -- (error) node[midway, above, label] {No};
        \draw[arrow] (error) -- (end1);
        \draw[arrow] (data) -- (template) node[midway, left, label] {Yes};
        \draw[arrow] (template) -- (pdf);
        \draw[arrow] (pdf) -- (download);
        \draw[arrow] (download) -- (end2);
    \end{tikzpicture}
    }
    \caption{Activity Diagram: Automated PDF Audit Report Generation}
    \label{fig:activity_report}
\end{figure}


\subsubsection{Framework Rule Mapping and Validation}
Administrators map Wazuh rules to framework controls as shown in Figure \ref{fig:activity_mapping}. Before saving the mapping, the system queries the SIEM API to verify the rule ID. If valid, the system checks for existing mappings in the database. This validation prevents accidental overwriting of existing configurations.

\begin{figure}[H]
    \centering
    \resizebox{0.8\textwidth}{!}{
    \begin{tikzpicture}[
        node distance=1.5cm,
        startstop/.style={circle, fill=black, minimum size=0.4cm},
        end/.style={circle, draw=black, thick, minimum size=0.4cm, path picture={\fill[black] (path picture bounding box.center) circle (0.15cm);}},
        activity/.style={rectangle, draw=black!80, thick, fill=blue!5, minimum width=3.5cm, minimum height=1cm, align=center, font=\scriptsize\bfseries, rounded corners=2pt, text width=3.2cm},
        decision/.style={diamond, draw=black!80, thick, fill=yellow!10, minimum width=2.5cm, minimum height=2.5cm, align=center, font=\scriptsize\bfseries, aspect=1.5, text width=1.8cm},
        arrow/.style={-Stealth, thick, draw=black!90},
        label/.style={font=\scriptsize, fill=white, inner sep=1pt}
    ]
        \node[startstop] (start) at (0,0) {};
        \node[activity] (init) at (0,-1.5) {Initiate Mapping Request};
        \node[activity] (api) at (0,-3) {Query Wazuh API for Rule ID};
        \node[decision] (exists) at (0,-5.5) {Rule Exists?};
        \node[activity] (error1) at (-4,-5.5) {Reject \& Show Error};
        \node[activity] (checkdb) at (0,-8.5) {Check DB for Existing Mapping};
        \node[decision] (collision) at (0,-11.5) {Collision Detected?};
        \node[activity] (override) at (4,-11.5) {Prompt Admin for Override};
        \node[activity] (commit) at (0,-14.5) {Commit Mapping to PostgreSQL};
        \node[activity] (audit) at (0,-16) {Write to Audit Trail};
        \node[end] (end) at (0,-17.5) {};

        \draw[arrow] (start) -- (init);
        \draw[arrow] (init) -- (api);
        \draw[arrow] (api) -- (exists);
        \draw[arrow] (exists) -- (error1) node[midway, above, label] {No};
        \draw[arrow] (exists) -- (checkdb) node[midway, right, label] {Yes};
        \draw[arrow] (checkdb) -- (collision);
        \draw[arrow] (collision) -- (override) node[midway, above, label] {Yes};
        \draw[arrow] (override) |- (commit);
        \draw[arrow] (collision) -- (commit) node[midway, right, label] {No};
        \draw[arrow] (commit) -- (audit);
        \draw[arrow] (audit) -- (end);
        \draw[arrow] (error1) |- (end);
    \end{tikzpicture}
    }
    \caption{Activity Diagram: Framework Rule Mapping and API Validation}
    \label{fig:activity_mapping}
\end{figure}

\subsubsection{Real-Time Alerting and Notification}
The platform includes an alerting pipeline as shown in Figure \ref{fig:activity_alerting}. When a compliance check fails, the engine evaluates its severity. Minor issues are recorded in the database logs, while critical failures trigger automated notifications via email and webhooks to auditors.

\begin{figure}[H]
    \centering
    \resizebox{0.8\textwidth}{!}{
    \begin{tikzpicture}[
        node distance=1.5cm,
        startstop/.style={circle, fill=black, minimum size=0.4cm},
        end/.style={circle, draw=black, thick, minimum size=0.4cm, path picture={\fill[black] (path picture bounding box.center) circle (0.15cm);}},
        activity/.style={rectangle, draw=black!80, thick, fill=blue!5, minimum width=3.5cm, minimum height=1cm, align=center, font=\scriptsize\bfseries, rounded corners=2pt, text width=3.2cm},
        decision/.style={diamond, draw=black!80, thick, fill=yellow!10, minimum width=2.5cm, minimum height=2.5cm, align=center, font=\scriptsize\bfseries, aspect=1.5, text width=2cm},
        fork/.style={rectangle, draw=black, fill=black, minimum width=6cm, minimum height=0.15cm},
        arrow/.style={-Stealth, thick, draw=black!90},
        label/.style={font=\scriptsize, fill=white, inner sep=1pt}
    ]
        \node[startstop] (start) at (0,0) {};
        \node[activity] (event) at (0,-1.5) {Receive State Change Event};
        \node[decision] (failed) at (0,-4) {Status == Failed?};
        \node[activity] (ignore) at (-5.5,-4) {Halt Execution};
        \node[decision] (severity) at (0,-7.5) {Severity $\ge$ Critical?};
        \node[activity] (log) at (-5.5,-7.5) {Write to Audit Log Only};
        \node[activity] (fetch) at (0,-10.5) {Fetch Subscribed Auditors};
        
        \node[fork] (fork1) at (0,-12) {};
        \node[activity] (email) at (-2.5,-13.5) {Dispatch Email via SMTP};
        \node[activity] (webhook) at (2.5,-13.5) {Trigger External Webhook};
        \node[fork] (join1) at (0,-15) {};
        
        \node[activity] (badge) at (0,-16.5) {Update UI Alert Badge};
        \node[end] (end) at (0,-18) {};

        \draw[arrow] (start) -- (event);
        \draw[arrow] (event) -- (failed);
        \draw[arrow] (failed) -- (ignore) node[midway, above, label] {No};
        \draw[arrow] (failed) -- (severity) node[midway, right, label] {Yes};
        \draw[arrow] (severity) -- (log) node[midway, above, label] {No};
        \draw[arrow] (severity) -- (fetch) node[midway, right, label] {Yes};
        \draw[arrow] (fetch) -- (fork1);
        \draw[arrow] (fork1) -| (email);
        \draw[arrow] (fork1) -| (webhook);
        \draw[arrow] (email) |- (join1);
        \draw[arrow] (webhook) |- (join1);
        \draw[arrow] (join1) -- (badge);
        \draw[arrow] (badge) -- (end);
        \draw[arrow] (ignore.west) -- ++(-1.5,0) |- (end);
        \draw[arrow] (log.west) -- ++(-1,0) |- (end);
    \end{tikzpicture}
    }
    \caption{Activity Diagram: Real-Time Alerting and Notification Fan-out}
    \label{fig:activity_alerting}
\end{figure}

\subsection{System Sequence Diagrams}
The sequence diagrams below trace the linear interactions between internal system components for each use case. Each diagram is identified by an SD-XX-Y code that maps directly to the parent UC-XX use case.

\subsubsection{SD-01: User Authentication}
The authentication sequence covers both the success and failure paths for UC-01. A valid credential submission results in JWT issuance, while an invalid submission terminates with an HTTP 401 response. These two paths are documented as SD-01-1 and SD-01-2 respectively.

Figure \ref{fig:sd_01_1} traces the success path. The client submits credentials via POST /login, the IAM module verifies the password hash against the database, signs a JWT, and returns it to the client.

\begin{figure}[H]
    \centering
    \resizebox{0.95\textwidth}{!}{
    \begin{tikzpicture}[
        node distance=2cm,
        inst/.style={rectangle, draw=black!90, thick, fill=blue!5, minimum width=2.8cm, minimum height=0.8cm, font=\small\bfseries, rounded corners=2pt},
        act/.style={rectangle, draw=black!90, fill=gray!20, minimum width=0.3cm},
        msg/.style={-Stealth, thick, draw=black!90},
        reply/.style={Stealth-, thick, dashed, draw=black!80},
        label/.style={font=\scriptsize, inner sep=2pt}
    ]
        \node[inst] (client) at (0, 0) {React SPA};
        \node[inst] (api) at (4.5, 0) {API Gateway};
        \node[inst] (iam) at (9, 0) {IAM Module};
        \node[inst] (db) at (13.5, 0) {PostgreSQL};
        \draw[dashed, draw=black!60, thick] (client.south) -- ++(0, -8);
        \draw[dashed, draw=black!60, thick] (api.south) -- ++(0, -8);
        \draw[dashed, draw=black!60, thick] (iam.south) -- ++(0, -8);
        \draw[dashed, draw=black!60, thick] (db.south) -- ++(0, -8);
        \draw[act] ([xshift=-0.15cm, yshift=-1cm]client.south) rectangle ([xshift=0.15cm, yshift=-7.5cm]client.south);
        \draw[act] ([xshift=-0.15cm, yshift=-1.2cm]api.south) rectangle ([xshift=0.15cm, yshift=-7.3cm]api.south);
        \draw[act] ([xshift=-0.15cm, yshift=-2cm]iam.south) rectangle ([xshift=0.15cm, yshift=-7cm]iam.south);
        \draw[act] ([xshift=-0.15cm, yshift=-3cm]db.south) rectangle ([xshift=0.15cm, yshift=-4.5cm]db.south);
        \draw[msg] ([yshift=-1.5cm]client.south) -- ([yshift=-1.5cm]api.south) node[midway, above=2pt, label] {1. POST /login};
        \draw[msg] ([yshift=-2.5cm]api.south) -- ([yshift=-2.5cm]iam.south) node[midway, above=2pt, label] {2. Verify()};
        \draw[msg] ([yshift=-3.5cm]iam.south) -- ([yshift=-3.5cm]db.south) node[midway, above=2pt, label] {3. Query User};
        \draw[reply] ([yshift=-4.2cm]iam.south) -- ([yshift=-4.2cm]db.south) node[midway, above=2pt, label] {4. Return Hash};
        \draw[msg] ([yshift=-5.2cm]iam.south) ++(0.15cm,0) -- ++(0.8,0) -- ++(0,-0.6) -- ++(-0.8,0) node[midway, right=4pt, label, text width=2cm] {5. Sign JWT};
        \draw[reply] ([yshift=-6.5cm]client.south) -- ([yshift=-6.5cm]api.south) node[midway, above=2pt, label] {6. 200 OK (JWT)};
    \end{tikzpicture}
    }
    \caption{SD-01-1: Valid Authentication Flow}
    \label{fig:sd_01_1}
\end{figure}

Figure \ref{fig:sd_01_2} traces the failure path. The submitted password hash does not match, and the system returns an HTTP 401 response without issuing any tokens.

\begin{figure}[H]
    \centering
    \resizebox{0.85\textwidth}{!}{
    \begin{tikzpicture}[
        node distance=2cm,
        inst/.style={rectangle, draw=black!90, thick, fill=blue!5, minimum width=2.8cm, minimum height=0.8cm, font=\small\bfseries, rounded corners=2pt},
        act/.style={rectangle, draw=black!90, fill=gray!20, minimum width=0.3cm},
        msg/.style={-Stealth, thick, draw=black!90},
        reply/.style={Stealth-, thick, dashed, draw=black!80},
        label/.style={font=\scriptsize, inner sep=2pt}
    ]
        \node[inst] (client) at (0, 0) {React SPA};
        \node[inst] (api) at (4.5, 0) {API Gateway};
        \node[inst] (iam) at (9, 0) {IAM Module};
        \draw[dashed, draw=black!60, thick] (client.south) -- ++(0, -6);
        \draw[dashed, draw=black!60, thick] (api.south) -- ++(0, -6);
        \draw[dashed, draw=black!60, thick] (iam.south) -- ++(0, -6);
        \draw[act] ([xshift=-0.15cm, yshift=-1cm]client.south) rectangle ([xshift=0.15cm, yshift=-5.5cm]client.south);
        \draw[act] ([xshift=-0.15cm, yshift=-1.2cm]api.south) rectangle ([xshift=0.15cm, yshift=-5.3cm]api.south);
        \draw[act] ([xshift=-0.15cm, yshift=-2cm]iam.south) rectangle ([xshift=0.15cm, yshift=-4cm]iam.south);
        \draw[msg] ([yshift=-1.5cm]client.south) -- ([yshift=-1.5cm]api.south) node[midway, above=2pt, label] {1. POST /login};
        \draw[msg] ([yshift=-2.5cm]api.south) -- ([yshift=-2.5cm]iam.south) node[midway, above=2pt, label] {2. Verify()};
        \draw[reply] ([yshift=-3.5cm]api.south) -- ([yshift=-3.5cm]iam.south) node[midway, above=2pt, label] {3. Hash Mismatch};
        \draw[reply] ([yshift=-4.8cm]client.south) -- ([yshift=-4.8cm]api.south) node[midway, above=2pt, label] {4. 401 Unauthorized};
    \end{tikzpicture}
    }
    \caption{SD-01-2: Invalid Authentication Flow}
    \label{fig:sd_01_2}
\end{figure}

\subsubsection{SD-02: Manage Wazuh Telemetry Synchronization}
This sequence traces UC-02, the telemetry synchronization pipeline. The Django API dispatches the task to a Celery worker, which polls the Wazuh SIEM, maps the returned rules to framework controls, and commits the resolved states to PostgreSQL.

Figure \ref{fig:sd_02_1} documents the linear success path for the synchronization operation.

\begin{figure}[H]
    \centering
    \resizebox{0.95\textwidth}{!}{
    \begin{tikzpicture}[
        node distance=2cm,
        inst/.style={rectangle, draw=black!90, thick, fill=blue!5, minimum width=2.8cm, minimum height=0.8cm, font=\small\bfseries, rounded corners=2pt},
        act/.style={rectangle, draw=black!90, fill=gray!20, minimum width=0.3cm},
        msg/.style={-Stealth, thick, draw=black!90},
        reply/.style={Stealth-, thick, dashed, draw=black!80},
        label/.style={font=\scriptsize, inner sep=2pt}
    ]
        \node[inst] (api) at (0, 0) {API Gateway};
        \node[inst] (celery) at (5, 0) {Celery Worker};
        \node[inst] (wazuh) at (10, 0) {Wazuh SIEM};
        \node[inst] (db) at (15, 0) {PostgreSQL};
        \draw[dashed, draw=black!60, thick] (api.south) -- ++(0, -8.5);
        \draw[dashed, draw=black!60, thick] (celery.south) -- ++(0, -8.5);
        \draw[dashed, draw=black!60, thick] (wazuh.south) -- ++(0, -8.5);
        \draw[dashed, draw=black!60, thick] (db.south) -- ++(0, -8.5);
        \draw[act] ([xshift=-0.15cm, yshift=-0.8cm]api.south) rectangle ([xshift=0.15cm, yshift=-2cm]api.south);
        \draw[act] ([xshift=-0.15cm, yshift=-1.5cm]celery.south) rectangle ([xshift=0.15cm, yshift=-8cm]celery.south);
        \draw[act] ([xshift=-0.15cm, yshift=-3cm]wazuh.south) rectangle ([xshift=0.15cm, yshift=-4.5cm]wazuh.south);
        \draw[act] ([xshift=-0.15cm, yshift=-6cm]db.south) rectangle ([xshift=0.15cm, yshift=-7.5cm]db.south);
        \draw[msg] ([yshift=-1.2cm]api.south) -- ([yshift=-1.2cm]celery.south) node[midway, above=2pt, label] {1. GET /sync};
        \draw[msg] ([yshift=-3.2cm]celery.south) -- ([yshift=-3.2cm]wazuh.south) node[midway, above=2pt, label] {2. GET /wazuh/agents};
        \draw[reply] ([yshift=-4.2cm]celery.south) -- ([yshift=-4.2cm]wazuh.south) node[midway, above=2pt, label] {3. 200 OK (JSON)};
        \draw[msg] ([yshift=-5.2cm]celery.south) ++(0.15cm,0) -- ++(0.8,0) -- ++(0,-0.6) -- ++(-0.8,0) node[midway, right=4pt, label, text width=2cm] {4. Map Rules()};
        \draw[msg] ([yshift=-6.5cm]celery.south) -- ([yshift=-6.5cm]db.south) node[midway, above=2pt, label] {5. Commit DB()};
        \draw[reply] ([yshift=-7.8cm]api.south) -- ([yshift=-7.8cm]celery.south) node[midway, above=2pt, label] {6. 200 OK};
    \end{tikzpicture}
    }
    \caption{SD-02-1: Telemetry Synchronization Success}
    \label{fig:sd_02_1}
\end{figure}

\subsubsection{SD-03: Map Framework Controls to Wazuh Rules}
This sequence covers UC-03, the framework-to-rule mapping operation. SD-03-1 documents the success path where a new mapping is validated and persisted. SD-03-2 documents the collision path where the backend detects and rejects a duplicate entry.

Figure \ref{fig:sd_03_1} traces the success path for creating a new control mapping.

\begin{figure}[H]
    \centering
    \resizebox{0.8\textwidth}{!}{
    \begin{tikzpicture}[
        node distance=2cm,
        inst/.style={rectangle, draw=black!90, thick, fill=blue!5, minimum width=2.8cm, minimum height=0.8cm, font=\small\bfseries, rounded corners=2pt},
        act/.style={rectangle, draw=black!90, fill=gray!20, minimum width=0.3cm},
        msg/.style={-Stealth, thick, draw=black!90},
        reply/.style={Stealth-, thick, dashed, draw=black!80},
        label/.style={font=\scriptsize, inner sep=2pt}
    ]
        \node[inst] (client) at (0, 0) {Admin UI};
        \node[inst] (api) at (4.5, 0) {API Gateway};
        \node[inst] (db) at (9, 0) {PostgreSQL};
        \draw[dashed, draw=black!60, thick] (client.south) -- ++(0, -6.5);
        \draw[dashed, draw=black!60, thick] (api.south) -- ++(0, -6.5);
        \draw[dashed, draw=black!60, thick] (db.south) -- ++(0, -6.5);
        \draw[act] ([xshift=-0.15cm, yshift=-1cm]client.south) rectangle ([xshift=0.15cm, yshift=-6cm]client.south);
        \draw[act] ([xshift=-0.15cm, yshift=-1.2cm]api.south) rectangle ([xshift=0.15cm, yshift=-5.5cm]api.south);
        \draw[act] ([xshift=-0.15cm, yshift=-2.5cm]db.south) rectangle ([xshift=0.15cm, yshift=-4.5cm]db.south);
        \draw[msg] ([yshift=-1.5cm]client.south) -- ([yshift=-1.5cm]api.south) node[midway, above=2pt, label] {1. POST /mappings};
        \draw[msg] ([yshift=-2.2cm]api.south) ++(0.15cm,0) -- ++(0.8,0) -- ++(0,-0.6) -- ++(-0.8,0) node[midway, right=4pt, label, text width=2cm] {2. Validate()};
        \draw[msg] ([yshift=-3.5cm]api.south) -- ([yshift=-3.5cm]db.south) node[midway, above=2pt, label] {3. Insert DB()};
        \draw[reply] ([yshift=-5.5cm]client.south) -- ([yshift=-5.5cm]api.south) node[midway, above=2pt, label] {4. 201 Created};
    \end{tikzpicture}
    }
    \caption{SD-03-1: Map Controls Success}
    \label{fig:sd_03_1}
\end{figure}

Figure \ref{fig:sd_03_2} traces the collision path where the backend detects a duplicate mapping and rejects it with HTTP 400.

\begin{figure}[H]
    \centering
    \resizebox{0.8\textwidth}{!}{
    \begin{tikzpicture}[
        node distance=2cm,
        inst/.style={rectangle, draw=black!90, thick, fill=blue!5, minimum width=2.8cm, minimum height=0.8cm, font=\small\bfseries, rounded corners=2pt},
        act/.style={rectangle, draw=black!90, fill=gray!20, minimum width=0.3cm},
        msg/.style={-Stealth, thick, draw=black!90},
        reply/.style={Stealth-, thick, dashed, draw=black!80},
        label/.style={font=\scriptsize, inner sep=2pt}
    ]
        \node[inst] (client) at (0, 0) {Admin UI};
        \node[inst] (api) at (4.5, 0) {API Gateway};
        \draw[dashed, draw=black!60, thick] (client.south) -- ++(0, -6);
        \draw[dashed, draw=black!60, thick] (api.south) -- ++(0, -6);
        \draw[act] ([xshift=-0.15cm, yshift=-1cm]client.south) rectangle ([xshift=0.15cm, yshift=-5.5cm]client.south);
        \draw[act] ([xshift=-0.15cm, yshift=-1.2cm]api.south) rectangle ([xshift=0.15cm, yshift=-5cm]api.south);
        \draw[msg] ([yshift=-1.5cm]client.south) -- ([yshift=-1.5cm]api.south) node[midway, above=2pt, label] {1. POST /mappings};
        \draw[msg] ([yshift=-2.5cm]api.south) ++(0.15cm,0) -- ++(0.8,0) -- ++(0,-0.6) -- ++(-0.8,0) node[midway, right=4pt, label, text width=2cm] {2. Validate()};
        \draw[reply] ([yshift=-3.8cm]api.south) ++(0.15cm,0) -- ++(0.8,0) -- ++(0,-0.6) -- ++(-0.8,0) node[midway, right=4pt, label, text width=2.5cm] {3. Duplicate Found()};
        \draw[reply] ([yshift=-5cm]client.south) -- ([yshift=-5cm]api.south) node[midway, above=2pt, label] {4. 400 Bad Request};
    \end{tikzpicture}
    }
    \caption{SD-03-2: Map Collision}
    \label{fig:sd_03_2}
\end{figure}

\subsubsection{SD-04: Monitor Endpoint Compliance Status}
This sequence traces UC-04, the dashboard compliance monitoring flow. The React SPA requests the aggregated compliance metrics from the API, which queries PostgreSQL and returns the calculated data array for visualization on the dashboard gauge.

Figure \ref{fig:sd_04_1} documents the data retrieval path for the compliance dashboard.

\begin{figure}[H]
    \centering
    \resizebox{0.85\textwidth}{!}{
    \begin{tikzpicture}[
        node distance=2cm,
        inst/.style={rectangle, draw=black!90, thick, fill=blue!5, minimum width=2.8cm, minimum height=0.8cm, font=\small\bfseries, rounded corners=2pt},
        act/.style={rectangle, draw=black!90, fill=gray!20, minimum width=0.3cm},
        msg/.style={-Stealth, thick, draw=black!90},
        reply/.style={Stealth-, thick, dashed, draw=black!80},
        label/.style={font=\scriptsize, inner sep=2pt}
    ]
        \node[inst] (client) at (0, 0) {React SPA};
        \node[inst] (api) at (5, 0) {API Gateway};
        \node[inst] (db) at (10, 0) {PostgreSQL};
        \draw[dashed, draw=black!60, thick] (client.south) -- ++(0, -7);
        \draw[dashed, draw=black!60, thick] (api.south) -- ++(0, -7);
        \draw[dashed, draw=black!60, thick] (db.south) -- ++(0, -7);
        \draw[act] ([xshift=-0.15cm, yshift=-1cm]client.south) rectangle ([xshift=0.15cm, yshift=-6.5cm]client.south);
        \draw[act] ([xshift=-0.15cm, yshift=-1.2cm]api.south) rectangle ([xshift=0.15cm, yshift=-6cm]api.south);
        \draw[act] ([xshift=-0.15cm, yshift=-2.5cm]db.south) rectangle ([xshift=0.15cm, yshift=-4cm]db.south);
        \draw[msg] ([yshift=-1.5cm]client.south) -- ([yshift=-1.5cm]api.south) node[midway, above=2pt, label] {1. GET /api/compliance};
        \draw[msg] ([yshift=-2.8cm]api.south) -- ([yshift=-2.8cm]db.south) node[midway, above=2pt, label] {2. Query Aggregated Metrics()};
        \draw[reply] ([yshift=-3.8cm]api.south) -- ([yshift=-3.8cm]db.south) node[midway, above=2pt, label] {3. Return Data Array};
        \draw[reply] ([yshift=-5.5cm]client.south) -- ([yshift=-5.5cm]api.south) node[midway, above=2pt, label] {4. 200 OK};
    \end{tikzpicture}
    }
    \caption{SD-04-1: Monitor Endpoint Compliance Status}
    \label{fig:sd_04_1}
\end{figure}

\subsubsection{SD-05: Resolve Compliance Anomalies}
This sequence traces UC-05, the anomaly remediation workflow. An administrator submits a remediation request through the UI. The API dispatches the task to a Celery worker, which updates the control status in PostgreSQL and confirms the resolution.

Figure \ref{fig:sd_05_1} documents the remediation dispatch and confirmation path.

\begin{figure}[H]
    \centering
    \resizebox{0.9\textwidth}{!}{
    \begin{tikzpicture}[
        node distance=2cm,
        inst/.style={rectangle, draw=black!90, thick, fill=blue!5, minimum width=2.8cm, minimum height=0.8cm, font=\small\bfseries, rounded corners=2pt},
        act/.style={rectangle, draw=black!90, fill=gray!20, minimum width=0.3cm},
        msg/.style={-Stealth, thick, draw=black!90},
        reply/.style={Stealth-, thick, dashed, draw=black!80},
        label/.style={font=\scriptsize, inner sep=2pt}
    ]
        \node[inst] (client) at (0, 0) {Admin UI};
        \node[inst] (api) at (4.5, 0) {API Gateway};
        \node[inst] (celery) at (9, 0) {Celery Worker};
        \node[inst] (db) at (13.5, 0) {PostgreSQL};
        \draw[dashed, draw=black!60, thick] (client.south) -- ++(0, -8);
        \draw[dashed, draw=black!60, thick] (api.south) -- ++(0, -8);
        \draw[dashed, draw=black!60, thick] (celery.south) -- ++(0, -8);
        \draw[dashed, draw=black!60, thick] (db.south) -- ++(0, -8);
        \draw[act] ([xshift=-0.15cm, yshift=-1cm]client.south) rectangle ([xshift=0.15cm, yshift=-7.5cm]client.south);
        \draw[act] ([xshift=-0.15cm, yshift=-1.2cm]api.south) rectangle ([xshift=0.15cm, yshift=-3cm]api.south);
        \draw[act] ([xshift=-0.15cm, yshift=-2.5cm]celery.south) rectangle ([xshift=0.15cm, yshift=-6.5cm]celery.south);
        \draw[act] ([xshift=-0.15cm, yshift=-4.5cm]db.south) rectangle ([xshift=0.15cm, yshift=-6cm]db.south);
        \draw[msg] ([yshift=-1.5cm]client.south) -- ([yshift=-1.5cm]api.south) node[midway, above=2pt, label] {1. POST /api/remediate};
        \draw[msg] ([yshift=-2.8cm]api.south) -- ([yshift=-2.8cm]celery.south) node[midway, above=2pt, label] {2. Dispatch Remediation Task()};
        \draw[msg] ([yshift=-5cm]celery.south) -- ([yshift=-5cm]db.south) node[midway, above=2pt, label] {3. Update Control Status()};
        \draw[reply] ([yshift=-5.8cm]celery.south) -- ([yshift=-5.8cm]db.south) node[midway, above=2pt, label] {4. Commit OK};
        \draw[reply] ([yshift=-7cm]client.south) -- ([yshift=-7cm]api.south) node[midway, above=2pt, label] {5. 200 OK};
    \end{tikzpicture}
    }
    \caption{SD-05-1: Resolve Compliance Anomalies}
    \label{fig:sd_05_1}
\end{figure}

\subsubsection{SD-06: Export Audit Reports}
This sequence traces UC-06, the PDF audit report generation pipeline. The auditor requests a compliance report, the API queries historical data from PostgreSQL, renders the PDF binary, and streams it back to the client browser.

Figure \ref{fig:sd_06_1} documents the report generation and delivery path.

\begin{figure}[H]
    \centering
    \resizebox{0.9\textwidth}{!}{
    \begin{tikzpicture}[
        node distance=2cm,
        inst/.style={rectangle, draw=black!90, thick, fill=blue!5, minimum width=2.8cm, minimum height=0.8cm, font=\small\bfseries, rounded corners=2pt},
        act/.style={rectangle, draw=black!90, fill=gray!20, minimum width=0.3cm},
        msg/.style={-Stealth, thick, draw=black!90},
        reply/.style={Stealth-, thick, dashed, draw=black!80},
        label/.style={font=\scriptsize, inner sep=2pt}
    ]
        \node[inst] (client) at (0, 0) {Auditor UI};
        \node[inst] (api) at (4.5, 0) {API Gateway};
        \node[inst] (db) at (9, 0) {PostgreSQL};
        \draw[dashed, draw=black!60, thick] (client.south) -- ++(0, -7);
        \draw[dashed, draw=black!60, thick] (api.south) -- ++(0, -7);
        \draw[dashed, draw=black!60, thick] (db.south) -- ++(0, -7);
        \draw[act] ([xshift=-0.15cm, yshift=-1cm]client.south) rectangle ([xshift=0.15cm, yshift=-6.5cm]client.south);
        \draw[act] ([xshift=-0.15cm, yshift=-1.2cm]api.south) rectangle ([xshift=0.15cm, yshift=-6cm]api.south);
        \draw[act] ([xshift=-0.15cm, yshift=-2.5cm]db.south) rectangle ([xshift=0.15cm, yshift=-3.5cm]db.south);
        \draw[msg] ([yshift=-1.5cm]client.south) -- ([yshift=-1.5cm]api.south) node[midway, above=2pt, label] {1. GET /reports/pdf};
        \draw[msg] ([yshift=-2.8cm]api.south) -- ([yshift=-2.8cm]db.south) node[midway, above=2pt, label] {2. Query Data()};
        \draw[msg] ([yshift=-4.2cm]api.south) ++(0.15cm,0) -- ++(0.8,0) -- ++(0,-0.6) -- ++(-0.8,0) node[midway, right=4pt, label, text width=2cm] {3. Render PDF()};
        \draw[reply] ([yshift=-5.8cm]client.south) -- ([yshift=-5.8cm]api.south) node[midway, above=2pt, label] {4. 200 OK (Binary Stream)};
    \end{tikzpicture}
    }
    \caption{SD-06-1: Report Generation}
    \label{fig:sd_06_1}
\end{figure}


\section{Low Level Design}
The Low-Level Design (LLD) translates the high-level system architecture into fundamental software building blocks. This section maps the specific object-oriented structures, persistence models, and relational bindings that dictate the execution of the ISCMS platform's core business logic.

\subsection{UML Class Diagram}
Figure \ref{fig:class_diagram} shows the Django ORM classes and their relationships. Relational constraints link regulatory frameworks to specific endpoint telemetry logs.

\begin{figure}[H]
    \centering
    \resizebox{\textwidth}{!}{
    \begin{tikzpicture}[
        classnode/.style={rectangle, draw=black!90, thick, fill=white, align=left, minimum width=4.5cm, rounded corners=2pt, inner sep=1ex, font=\scriptsize},
        assoc/.style={-Stealth, thick, draw=black!80},
        rel/.style={font=\scriptsize\bfseries, fill=white, inner sep=1pt}
    ]
        % Classes
        \node[classnode] (user) at (0, 0) {
            \textbf{UserAccount}\\[1ex]
            \rule{4cm}{0.4pt}\\[0.5ex]
            + uuid : UUID\\
            + email : String\\
            + password\_hash : String\\
            + role : Enum (Admin, Auditor)\\
            + is\_active : Boolean\\[0.5ex]
            \rule{4cm}{0.4pt}\\[0.5ex]
            + authenticate() : Boolean\\
            + generate\_token\_pair() : JWT
        };
        \node[classnode] (framework) at (7, 0) {
            \textbf{RegulatoryFramework}\\[1ex]
            \rule{4cm}{0.4pt}\\[0.5ex]
            + id : Integer\\
            + name : String (e.g., ISO 27001)\\
            + version : String\\
            + description : Text\\[0.5ex]
            \rule{4cm}{0.4pt}\\[0.5ex]
            + get\_controls() : QuerySet\\
            + calculate\_compliance() : Float
        };
        \node[classnode] (control) at (14, 0) {
            \textbf{FrameworkControl}\\[1ex]
            \rule{4cm}{0.4pt}\\[0.5ex]
            + id : Integer\\
            + framework\_id : FK\\
            + clause\_id : String\\
            + objective : Text\\[0.5ex]
            \rule{4cm}{0.4pt}\\[0.5ex]
            + evaluate\_status() : Boolean
        };
        \node[classnode] (mapping) at (14, -6) {
            \textbf{RuleMapping}\\[1ex]
            \rule{4cm}{0.4pt}\\[0.5ex]
            + id : Integer\\
            + control\_id : FK\\
            + wazuh\_rule\_id : Integer\\
            + severity\_weight : Integer\\[0.5ex]
            \rule{4cm}{0.4pt}\\[0.5ex]
            + validate\_mapping() : Boolean
        };
        \node[classnode] (endpoint) at (7, -6) {
            \textbf{MonitoredEndpoint}\\[1ex]
            \rule{4cm}{0.4pt}\\[0.5ex]
            + agent\_id : String\\
            + hostname : String\\
            + ip\_address : String\\
            + status : Enum\\[0.5ex]
            \rule{4cm}{0.4pt}\\[0.5ex]
            + fetch\_latest\_telemetry() : JSON
        };
        \node[classnode] (state) at (0, -6) {
            \textbf{ComplianceState}\\[1ex]
            \rule{4cm}{0.4pt}\\[0.5ex]
            + id : Integer\\
            + endpoint\_id : FK\\
            + rule\_id : FK\\
            + is\_passed : Boolean\\
            + timestamp : DateTime\\[0.5ex]
            \rule{4cm}{0.4pt}\\[0.5ex]
            + trigger\_alert() : Void\\
            + commit\_to\_ledger() : Void
        };
        % Relationships
        \draw[assoc] (framework.east) -- (control.west) node[midway, above, rel] {1} node[midway, below, rel] {1..*};
        \draw[assoc] (control.south) -- (mapping.north) node[midway, left, rel] {1} node[midway, right, rel] {1..*};
        \draw[assoc] (endpoint.west) -- (state.east) node[midway, above, rel] {1} node[midway, below, rel] {1..*};
        \draw[assoc, dashed] (mapping.west) -- (endpoint.east) node[midway, above, rel] {Evaluates};
        
    \end{tikzpicture}
    }
    \caption{UML Class Diagram: Core Object-Relational Models}
    \label{fig:class_diagram}
\end{figure}

\section{Database Design}
I built the database on PostgreSQL. The tables are normalized to the Third Normal Form (3NF) to prevent update anomalies. When you're pulling in thousands of SIEM alerts an hour, you need strict ACID compliance so the data doesn't corrupt.

\subsection{Physical Data Model (Core ERD)}
Figure \ref{fig:core_erd} shows the actual physical schema. I mapped out the exact column types and foreign keys we used to build the PostgreSQL database.

\begin{figure}[H]
    \centering
    \resizebox{\textwidth}{!}{
    \begin{tikzpicture}[
        table/.style={rectangle, draw=black!90, thick, fill=white, align=left, font=\footnotesize\ttfamily, inner sep=1.5ex, rounded corners=2pt},
        rel/.style={-Stealth, thick, draw=black!80},
        card/.style={font=\scriptsize\bfseries\sffamily, fill=white, inner sep=1pt}
    ]

        % ----------------------------------------------------
        % TABLE PLACEMENT (Strict Grid Layout)
        % ----------------------------------------------------
        \node[table, text width=4.5cm] (users) at (0, 6) {
            \textbf{users} \\
            \rule{4.5cm}{0.4pt} \\
            PK \hspace{1mm} id : UUID \\
            UQ \hspace{1mm} email : VARCHAR(255) \\
            \hspace{5.5mm} password\_hash : VARCHAR(255) \\
            \hspace{5.5mm} role : VARCHAR(50) \\
            \hspace{5.5mm} created\_at : TIMESTAMP
        };

        \node[table, text width=4.5cm] (frameworks) at (8, 6) {
            \textbf{regulatory\_frameworks} \\
            \rule{4.5cm}{0.4pt} \\
            PK \hspace{1mm} id : SERIAL \\
            UQ \hspace{1mm} name : VARCHAR(255) \\
            \hspace{5.5mm} version : VARCHAR(50) \\
            \hspace{5.5mm} is\_active : BOOLEAN
        };

        \node[table, text width=4.5cm] (controls) at (16, 6) {
            \textbf{framework\_controls} \\
            \rule{4.5cm}{0.4pt} \\
            PK \hspace{1mm} id : SERIAL \\
            FK \hspace{1mm} framework\_id : INT \\
            \hspace{5.5mm} clause\_id : VARCHAR(50) \\
            \hspace{5.5mm} objective : TEXT
        };

        \node[table, text width=4.5cm] (endpoints) at (0, 0) {
            \textbf{monitored\_endpoints} \\
            \rule{4.5cm}{0.4pt} \\
            PK \hspace{1mm} agent\_id : VARCHAR(50) \\
            FK \hspace{1mm} managed\_by : UUID \\
            \hspace{5.5mm} hostname : VARCHAR(255) \\
            \hspace{5.5mm} ip\_address : INET \\
            \hspace{5.5mm} status : VARCHAR(20)
        };

        \node[table, text width=4.5cm] (states) at (8, 0) {
            \textbf{compliance\_states} \\
            \rule{4.5cm}{0.4pt} \\
            PK \hspace{1mm} id : BIGSERIAL \\
            FK \hspace{1mm} agent\_id : VARCHAR(50) \\
            FK \hspace{1mm} rule\_mapping\_id : INT \\
            \hspace{5.5mm} status : VARCHAR(20) \\
            \hspace{5.5mm} timestamp : TIMESTAMP
        };

        \node[table, text width=4.5cm] (mappings) at (16, 0) {
            \textbf{rule\_mappings} \\
            \rule{4.5cm}{0.4pt} \\
            PK \hspace{1mm} id : SERIAL \\
            FK \hspace{1mm} control\_id : INT \\
            UQ \hspace{1mm} wazuh\_rule\_id : INT \\
            \hspace{5.5mm} severity\_weight : INT
        };

        % ----------------------------------------------------
        % RELATIONSHIPS (Crow's Foot via 1-to-N lines)
        % ----------------------------------------------------
        \draw[rel] (users.south) -- (endpoints.north) node[pos=0.1, right, card] {1} node[pos=0.9, right, card] {N};
        \draw[rel] (frameworks.east) -- (controls.west) node[pos=0.1, above, card] {1} node[pos=0.9, above, card] {N};
        \draw[rel] (controls.south) -- (mappings.north) node[pos=0.1, right, card] {1} node[pos=0.9, right, card] {N};
        \draw[rel] (endpoints.east) -- (states.west) node[pos=0.1, above, card] {1} node[pos=0.9, above, card] {N};
        \draw[rel] (mappings.west) -- (states.east) node[pos=0.1, above, card] {1} node[pos=0.9, above, card] {N};

    \end{tikzpicture}
    }
    \caption{Physical Database Schema: Core Relational Data Model with Exact Data Types}
    \label{fig:core_erd}
\end{figure}

\subsection{Extended Entity Relationship Diagram (EERD)}
Figure \ref{fig:eerd} is an Extended ERD (EERD) that visualizes the logical rules we put in place.

Notice the disjoint (\textbf{d}) specialization: a generic User Account splits strictly into either a System Admin or a Security Auditor. That handles our RBAC at the database level. I also modeled the Compliance State as a weak entity, because an alert record is useless without its parent Monitored Endpoint.

\begin{figure}[H]
    \centering
    \resizebox{0.9\textwidth}{!}{
    \begin{tikzpicture}[
        entity/.style={rectangle, draw=black!90, thick, fill=blue!10, minimum width=2.8cm, minimum height=1cm, align=center, font=\small\bfseries, rounded corners=2pt},
        weakent/.style={rectangle, double, draw=black!90, thick, fill=blue!5, minimum width=2.8cm, minimum height=1cm, align=center, font=\small\bfseries, rounded corners=2pt},
        rel/.style={diamond, draw=black!90, thick, fill=yellow!20, minimum width=2.5cm, minimum height=1.5cm, align=center, font=\footnotesize\bfseries, aspect=1.5},
        weakrel/.style={diamond, double, draw=black!90, thick, fill=yellow!10, minimum width=2.5cm, minimum height=1.5cm, align=center, font=\footnotesize\bfseries, aspect=1.5},
        disjoint/.style={circle, draw=black!90, thick, fill=white, minimum size=0.6cm, font=\small\bfseries},
        line/.style={thick, draw=black!80},
        attr/.style={ellipse, draw=black!80, fill=white, font=\footnotesize, inner sep=1.5pt},
        card/.style={font=\scriptsize\bfseries, fill=white, inner sep=1pt}
    ]

        % ----------------------------------------------------
        % CORE ENTITIES (Center Y-Axis Stack)
        % ----------------------------------------------------
        \node[entity] (user) at (0, 5) {User Account};
        \node[disjoint] (d) at (0, 3) {d};
        \node[font=\scriptsize\bfseries, right=2pt] at (0, 4) {Is-A};
        
        \node[entity] (admin) at (-4, 1) {System Admin};
        \node[entity] (auditor) at (4, 1) {Security Auditor};

        \node[rel] (monitors) at (-4, -1.5) {Monitors};
        \node[rel] (reviews) at (4, -1.5) {Reviews};

        \node[entity] (endpoint) at (-4, -4) {Monitored Endpoint};
        \node[weakrel] (generates) at (0, -4) {Generates};
        \node[weakent] (state) at (4, -4) {Compliance State};

        % ----------------------------------------------------
        % ATTRIBUTES (Shooting Outwards to Prevent Overlap)
        % ----------------------------------------------------
        % User Attributes (Top)
        \node[attr] (uuid) at (0, 7) {\underline{uuid}};
        \node[attr] (email) at (-2.5, 6) {email};
        \node[attr] (pass) at (2.5, 6) {password};
        \draw[line] (user.north) -- (uuid.south);
        \draw[line] (user.north west) -- (email.south east);
        \draw[line] (user.north east) -- (pass.south west);

        % Subclass Attributes (Far Left / Far Right)
        \node[attr] (clearance) at (-8.0, 1) {clearance\_level};
        \draw[line] (admin.west) -- (clearance.east);

        \node[attr] (cert) at (8.0, 1) {cert\_id};
        \draw[line] (auditor.east) -- (cert.west);

        % Endpoint Attributes (Bottom Left)
        \node[attr] (agent) at (-8.0, -4) {\underline{agent\_id}};
        \node[attr] (hostname) at (-6.5, -5.5) {hostname};
        \node[attr] (ip) at (-4, -6.5) {ip\_address};
        \draw[line] (endpoint.west) -- (agent.east);
        \draw[line] (endpoint.south west) -- (hostname.north east);
        \draw[line] (endpoint.south) -- (ip.north);

        % State Attributes (Bottom Right)
        \node[attr] (timestamp) at (8.5, -4) {\textit{timestamp (Partial)}};
        \node[attr] (status) at (7.5, -5.5) {status};
        \node[attr] (details) at (4, -6.5) {details};
        \draw[line] (state.east) -- (timestamp.west);
        \draw[line] (state.south east) -- (status.north west);
        \draw[line] (state.south) -- (details.north);

        % ----------------------------------------------------
        % STRUCTURAL CONNECTIONS
        % ----------------------------------------------------
        % Inheritance
        \draw[line] (user.south) -- (d.north);
        \draw[line] (d.south) -- ++(0,-0.5) -| (admin.north);
        \draw[line] (d.south) -- ++(0,-0.5) -| (auditor.north);

        % Admin to Endpoint
        \draw[line] (admin.south) -- (monitors.north) node[midway, left, card] {1};
        \draw[line] (monitors.south) -- (endpoint.north) node[midway, left, card] {N};

        % Auditor to State
        \draw[line] (auditor.south) -- (reviews.north) node[midway, right, card] {1};
        \draw[line] (reviews.south) -- (state.north) node[midway, right, card] {N};

        % Weak Entity Dependency
        \draw[line] (endpoint.east) -- (generates.west) node[midway, above, card] {1};
        \draw[line, double] (generates.east) -- (state.west) node[midway, above, card] {N};

    \end{tikzpicture}
    }
    \caption{Extended Entity Relationship Diagram (EERD) with Constraints and Attributes}
    \label{fig:eerd}
\end{figure}


\section{Graphical User Interface (GUI) Design}
The frontend is a React Single Page Application (SPA) designed to present compliance telemetry without visual clutter. The screenshots below demonstrate the default dark mode, though the application supports dynamic theme switching.

\subsection{Authentication and Identity Management}
The authentication portal validates user credentials and issues a JSON Web Token (JWT). The application then initializes the dashboard corresponding to the user's RBAC role.

\begin{figure}[H]
    \centering
    \includegraphics[width=1.0\textwidth]{"images/login.png"}
    \caption{GUI Design: Secure User Authentication Portal.}
    \label{fig:gui_login}
\end{figure}

\begin{figure}[H]
    \centering
    \includegraphics[width=1.0\textwidth]{"images/2fa.png"}
    \caption{GUI Design: Time-based One-Time Password (TOTP) Verification.}
    \label{fig:gui_2fa}
\end{figure}

\subsection{Superadmin Configuration Portal}
Superadmins have full administrative privileges, including configuration of the Wazuh pipeline, regulatory mapping rules, and user account provisioning. The interface supports dynamic theming for varied lighting conditions.

\begin{figure}[H]
    \centering
    \includegraphics[width=1.0\textwidth]{"images/superadmin dashboard dark.png"}
    \caption{GUI Design: Superadmin Global Compliance Dashboard (Dark Mode).}
    \label{fig:gui_superadmin_dark}
\end{figure}

\begin{figure}[H]
    \centering
    \includegraphics[width=1.0\textwidth]{"images/superadmin dashboard1.png"}
    \caption{GUI Design: Superadmin Global Compliance Dashboard (Light Mode).}
    \label{fig:gui_superadmin_light}
\end{figure}

\begin{figure}[H]
    \centering
    \includegraphics[width=1.0\textwidth]{"images/manage admins1.png"}
    \caption{GUI Design: RBAC User and System Administrator Management Interface.}
    \label{fig:gui_manage_admins}
\end{figure}

\subsection{Auditor Workspace (RBAC Enforcement)}
The application modifies the render tree based on the authenticated role. Auditors are restricted to a read-only interface to prevent unauthorized modifications to audit evidence or framework rules. Configuration and management modules are omitted from this view.

\begin{figure}[H]
    \centering
    \includegraphics[width=1.0\textwidth]{"images/auditor dashboard.png"}
    \caption{GUI Design: Security Auditor Read-Only Compliance View.}
    \label{fig:gui_auditor_dashboard}
\end{figure}

\subsection{Automated Audit Artifacts and Notifications}
The system provides PDF audit reports for external reviewers and SMTP alerts for administrators. The backend service aggregates PostgreSQL data to generate these artifacts and sends notifications when compliance thresholds are breached.

\begin{figure}[H]
    \centering
    \includegraphics[width=1.0\textwidth]{"images/compliance report pdf.png"}
    \caption{System Output: Automated PDF Audit Dossier Generation.}
    \label{fig:gui_pdf_report}
\end{figure}

\begin{figure}[H]
    \centering
    \includegraphics[width=1.0\textwidth]{"images/email.png"}
    \caption{System Output: SMTP Alert for Rule Violation.}
    \label{fig:gui_email_alert}
\end{figure}


\section{External Interfaces}
The ISCMS doesn't run in a vacuum. It talks to several external services. This section lists out the protocols and authentication methods I used to keep those connections secure.

\subsection{Wazuh SIEM REST API}
The whole system relies on raw logs coming from the Wazuh Manager. My Celery workers hit the Wazuh REST API over HTTPS. I secured this with Bearer tokens generated by the Wazuh server. We pull down the JSON payloads and validate them against the Django ORM before saving them.

\subsection{SMTP Notification Gateway}
For real-time alerts, the backend constructs HTML emails and fires them off through an external SMTP server over port 587 using STARTTLS. I threw this logic into the async message broker so it wouldn't block the API Gateway.

\subsection{Outbound Webhook Integration}
I also added outbound webhooks so the system can talk to tools like Slack or Microsoft Teams. Whenever a compliance rule breaks, it shoots out an HTTP POST request with the JSON event data. To prevent spoofing, I added an HMAC-SHA256 signature to the headers.

\subsection{Interface Specification Matrix}
Table \ref{tab:external_interfaces} summarizes the technical parameters and boundary constraints of the platform's external interfaces.

\begin{table}[H]
    \centering
    \caption{External Interface Technical Specifications}
    \label{tab:external_interfaces}
    \renewcommand{\arraystretch}{1.4}
    \resizebox{0.95\textwidth}{!}{
    \begin{tabular}{|l|l|l|l|l|}
        \hline
        \rowcolor{blue!10} 
        \textbf{Interface Target} & \textbf{Network Protocol} & \textbf{Data Format} & \textbf{Authentication Mechanism} & \textbf{Directionality} \\ \hline
        Wazuh Manager API & HTTPS (TCP 443) & JSON & Bearer Token / API Key & Bidirectional (Polling) \\ \hline
        SMTP Relay Server & SMTP (TCP 587) & MIME / HTML & SMTP AUTH (TLS) & Outbound Only \\ \hline
        Enterprise Webhooks & HTTPS (TCP 443) & JSON & HMAC-SHA256 Signatures & Outbound Only \\ \hline
        Client Web Browser & HTTPS (TCP 443) & HTML / JSON & JWT (Stateless) & Bidirectional \\ \hline
    \end{tabular}
    }
\end{table}
